% บทคัดย่อภาษาไทย
โครงงานฉบับนี้มีวัตถุประสงค์เพื่อออกแบบและพัฒนาเกมเดินเรือด้วยระบบสั่งการด้วยเสียง (Voice-Controlled Ship Navigation Game) สำหรับเป็นสื่อการเรียนรู้เพื่อส่งเสริมทักษะการสื่อสารและลดความวิตกกังวลในการใช้ภาษาต่างประเทศของผู้เรียน โดยระบบผสมผสานรูปแบบเกมจำลองสถานการณ์เข้ากับเทคโนโลยีการรู้จำเสียง (Speech Recognition) ประกอบด้วยระบบหลัก ได้แก่ ระบบรู้จำเสียงและประมวลผลคำสั่ง ระบบฟิสิกส์และการควบคุมเรือ ระบบลูกเรือปัญญาประดิษฐ์ ระบบสร้างสิ่งกีดขวาง ระบบส่วนติดต่อผู้ใช้ ระบบช่วยสอนภาษา โหมดเนื้อเรื่อง และระบบบันทึกผลการทดลอง ผู้เล่นรับบทเป็นกัปตันเรือที่ต้องสื่อสารและสั่งการลูกเรือปัญญาประดิษฐ์ ได้แก่ ต้นกลเรือผู้ควบคุมเครื่องยนต์ และพลถือท้ายและยามระวังหน้าผู้บังคับหางเสือและเฝ้าระวังสิ่งกีดขวาง ด้วยวลีสื่อสารมาตรฐานทางทะเล ระบบพัฒนาด้วยเอนจิน Unity และภาษา C\# ระบบแปลงคำสั่งเสียงภาษาอังกฤษเป็นข้อความภายในเครื่องของผู้เล่นด้วยโมเดล Whisper ผ่านไลบรารี whisper.cpp แล้วส่งให้ตัวแยกคำสั่งตีความ หากตัวแยกคำสั่งไม่เข้าใจคำสั่ง ระบบจึงใช้ GPT-4o mini ช่วยตีความ ระบบลูกเรือปัญญาประดิษฐ์จำลองสถานการณ์สื่อสารขัดข้องเมื่อผู้เล่นออกเสียงไม่ชัดเจนหรือใช้คำสั่งผิดบริบท โดยลูกเรือจะทวนถามหรือขอให้พูดซ้ำ เกมจึงออกแบบมาเพื่อให้ผู้เล่นได้ฝึกความคล่องแคล่วในการพูดและการแก้ปัญหาเฉพาะหน้าภายใต้ความกดดัน ในรายงานความก้าวหน้าครั้งนี้ ผู้จัดทำได้พัฒนาต้นแบบของเกมแล้ว ผ่านการทดสอบหน่วย (Unit Test) ครบทั้ง 363 จาก 363 กรณี และการทดสอบระบบโดยผู้จัดทำ ส่วนการประเมินกับผู้ใช้จะดำเนินการในระยะถัดไป โดยจะประเมินประสิทธิภาพของระบบจากแฟ้มบันทึกผลภายในเกม และเปรียบเทียบความวิตกกังวลในการพูดภาษาอังกฤษของกลุ่มตัวอย่างจำนวน 10--20 คนก่อนและหลังเล่นเกมด้วยแบบวัด 12 ข้อที่ดัดแปลงจาก Foreign Language Classroom Anxiety Scale (FLCAS) ร่วมกับแบบสอบถามประสบการณ์การเล่นเกม
